%%%%%%%%%%%%%%%%%%%%%%%%%%%%%%%%%%%%%%%%%%%%%%%%%%%%%%%%%%%%%%%%%%%%%%%%%%%%%%%%
% AgentShield -- submission manuscript
%
% Target venue : USENIX Security 2027 (33rd USENIX Security Symposium),
%                Cycle 2.
%                Registration deadline 2027-01-19, paper deadline 2027-01-26,
%                artifact submission 2027-01-29.
% Source       : docs/research/MANUSCRIPT_v041_en.md (frozen internal research
%                release; the release tag and all revision identifiers are
%                intentionally absent from this file and from the compiled PDF,
%                because USENIX Security double-blind review forbids
%                identifying repository history)
% Figures      : paper/figures/*.pdf, vector output of
%                scripts/draw_manuscript_figures.py
%
% Build        : pdflatex main && bibtex main && pdflatex main && pdflatex main
%
% Anonymous    : no author name, affiliation, funding, acknowledgement, public
%                repository owner/name, git tag, git commit SHA, CI run link or
%                local filesystem path may appear in any form, including inside
%                comments, PDF metadata and embedded file names
%
% Limits       : 13-page body (references and appendices are excluded);
%                unlimited appendices at submission time; 20 pages camera-ready.
%                Open Science appendix is mandatory, Ethics appendix is
%                strongly encouraged.  Anonymous submission: no acknowledgements
%                that identify the authors.
%%%%%%%%%%%%%%%%%%%%%%%%%%%%%%%%%%%%%%%%%%%%%%%%%%%%%%%%%%%%%%%%%%%%%%%%%%%%%%%%

\documentclass[letterpaper,twocolumn,10pt]{article}
\usepackage{usenix}

\usepackage{amsmath}
\usepackage{booktabs}
\usepackage{graphicx}

\graphicspath{{figures/}}

% Keeps the compiled PDF's document properties free of identifying metadata:
% usenix.sty loads hyperref, so these are the values the PDF will carry.
\hypersetup{
  pdftitle={AgentShield: Closing the Observability Gap in Tool-Using LLM Agents
            with Provenance-Aware Runtime Governance},
  pdfauthor={Anonymous Submission},
  pdfcreator={},
  pdfproducer={},
  pdfsubject={USENIX Security 2027 submission (double-blind)}
}

% Lets long identifiers and hex digests break across lines inside \texttt.
\newcommand{\hb}{\allowbreak}

%-------------------------------------------------------------------------------
\begin{document}
%-------------------------------------------------------------------------------

\date{}

% make title bold and 14 pt font (Latex default is non-bold, 16 pt)
\title{\Large \bf AgentShield: Closing the Observability Gap in Tool-Using
LLM Agents with Provenance-Aware Runtime Governance}

% anonymous submission block: no names, no institutions, no funding
\author{
{\rm Anonymous Authors}\\
Anonymous Submission
}

\maketitle

%-------------------------------------------------------------------------------
\begin{abstract}
%-------------------------------------------------------------------------------
Tool-using large language model (LLM) agents increasingly act on information
obtained from external tools, files, web pages, and other agents. This creates
an observability problem for runtime guardrails: the decisive malicious
instruction may occur in an earlier tool output, while the eventual tool call
is locally ordinary. We study this gap and present AgentShield, a
provenance-aware runtime governance prototype that records tool outputs as
typed artifacts, tracks the first origin of entities such as destinations and
file paths, propagates risk across a behavior graph, and gates each tool call
with ALLOW, HUMAN\_REVIEW, or BLOCK. The evaluation is deliberately separated
from grading metadata by an explicit evaluation contract and a
prediction-invariance test over 3,360 metadata permutations, under which no
prediction changes. On an AgentDojo-derived logged-trace replay containing all
97 benign trajectories and a deterministic prefix of 400 attack-labelled
trajectories, a current-call-only gate
blocks 0\% of attack trajectories. Inspecting untrusted prior tool output
raises the attack block rate to 3.5\%, and adding entity provenance raises it
to 16.0\%. The same progression increases benign trace blocking from 1.0\% to
19.6\% and 41.2\%, respectively, exposing a substantial safety-utility
trade-off rather than hiding it.

We then show that part of the observed benign blocking is attributable to
authorization and risk-aggregation semantics rather than to provenance alone.
When operator authorization is allowed to temper only provenance-derived
alarms and is structurally prevented from pardoning independently dangerous
behavior, benign trace blocking falls to
28.9\% while the raw blocking rate on attack-labelled trajectories falls to
3.5\%. An action-level audit finds that the released final calls are
predominantly authorization-concordant or read-only; offline replay cannot
determine whether the underlying attacks would still succeed. A paired
counterfactual test further shows that identical final tool calls can receive
different decisions when their source context differs. However, a frozen
string-based instruction detector achieves 0/9 recall on differently phrased
held-out injections, and corpus-specific markers remain disabled by default.
These results support provenance as a useful missing runtime signal while
identifying trust attribution and instruction-data separation as the primary
unresolved problems.
\end{abstract}

\noindent{\small\textbf{Keywords:} LLM agents; prompt injection; provenance;
taint tracking; runtime governance; behavior graphs; agent security;
authorization semantics}

%-------------------------------------------------------------------------------
\section{Introduction}
\label{sec:intro}
%-------------------------------------------------------------------------------

LLM agents differ from chat systems in a security-critical way: they do not
merely generate text, but select tools and cause side effects. An agent can
read a calendar, inspect a document, query a database, call an API, transfer
funds, send an email, or delete a file. At the same time, many of the inputs an
agent consumes are not authored by the operator. They arrive through tool
responses, web pages, retrieved documents, email bodies, or other agents.
Indirect prompt injection exploits exactly this boundary by placing adversarial
instructions in data that later becomes part of the model
context~\cite{zhan2024injecagent,debenedetti2024agentdojo,chen2025struq}.

Most runtime guardrails are naturally attached to the point where an agent is
about to execute a tool. This is attractive because the enforcement point is
concrete: inspect the tool name and arguments, estimate risk, then allow,
review, or block. Yet this placement can create a structural observability
gap. Suppose the operator asks an agent to summarize a document. The document
contains an injected instruction to send a copy to an attacker-controlled
address. When the model later calls \texttt{send\_email}, the current call may
contain a syntactically ordinary recipient and subject. The information that
makes the call suspicious is not necessarily present in the call itself; it is
the fact that the recipient was introduced by an untrusted artifact rather than
by the operator.

This paper asks a narrower question than end-to-end prompt-injection defense:
can runtime governance improve by tracking the provenance of values that flow
from prior observations into later actions? AgentShield is a research
prototype designed to study that question. It turns observed tool outputs into
first-class artifacts, extracts entities from those artifacts, records the
first origin and trust level of each entity, links tool calls in a behavior
graph, and computes deterministic provenance signals before each
side-effecting call.

Our contribution is not a claim of solved prompt injection. The frozen
evaluation shows the opposite. Provenance materially increases attack
interception in logged traces, but under a deliberately conservative
all-untrusted trust policy it also increases benign blocking. Moreover, a
string-marker instruction detector fails completely on a small set of
differently phrased held-out attacks. We therefore frame AgentShield as
evidence for a missing signal and as an instrument for measuring the
trade-off, not as a complete defense.

The v0.4.1 configuration reported here is a \emph{correctness} release rather
than a performance release. It changes how authorization interacts with risk
aggregation, and the measured cost of that change is reported rather than
optimized away. No threshold, weight, marker, or sampling rule was retuned to
recover any number.

The paper makes five contributions:

\begin{itemize}

\item A runtime formulation of the observability gap: a tool call can be
locally benign-looking even when its arguments originate in adversarial
earlier context.

\item A provenance-aware governance mechanism that records content artifacts,
tracks entity origin and taint, propagates risk through a behavior graph, and
applies a two-pass decision after propagation.

\item \textbf{An authorization safety invariant for risk aggregation}
(Section~\ref{sec:invariant}): operator authorization may explain
provenance-derived suspicion, but it cannot authorize away independently
dangerous behavior. We show that the natural implementation --- a single cap on
the combined score --- is an authorization bypass in which an authorized bulk
delete is scored as a benign action, and we give a band-split formulation that
closes it.

\item An evaluation contract that physically separates runtime-observable
inputs from evaluation-only metadata, together with a prediction-invariance
test over 3,360 metadata perturbations.

\item An evaluation that reports both the detection gain and its benign cost,
including negative generalization results and explicit non-measurement of
attack success rate, benign utility, and task success.

\end{itemize}

%-------------------------------------------------------------------------------
\section{Related Work}
\label{sec:related}
%-------------------------------------------------------------------------------

\subsection{Indirect prompt injection in tool-integrated agents}
\label{sec:rw:ipi}

InjecAgent demonstrates that indirect prompt injection can manipulate
tool-integrated agents through externally supplied content, and evaluates a
large collection of injected agent tasks~\cite{zhan2024injecagent}. AgentDojo
provides a dynamic environment for evaluating attacks and defenses in
tool-using agents and makes explicit that untrusted tool outputs can carry
instructions that conflict with the operator
intent~\cite{debenedetti2024agentdojo}. These benchmarks motivate the threat
model studied here: dangerous behavior may emerge only after the model has
consumed earlier external content.

Subsequent work broadened the attack surface in directions our replay does not
model. ToolHijacker poisons the tool document itself rather than a tool
\emph{response}, manipulating tool
selection~\cite{shi2025toolselection}. Automated optimization-based attackers
inject adversarial suffixes learned by reinforcement learning and transfer them
to models fine-tuned for injection
resistance~\cite{chen2026learningtoinject}. Role confusion offers a mechanistic
account, showing that injected text occupies the same representational space as
the role it imitates~\cite{ye2026roleconfusion}. These results bound the scope
of our contribution: we study the runtime governance point, not the model's
internal role perception, and our measurements are taken against the frozen
dump rather than against adaptive attackers.

\subsection{Input separation and provenance-aware defenses}
\label{sec:rw:separation}

StruQ separates instructions from data through structured queries and trains
models to follow only the instruction
channel~\cite{chen2025struq}. Spotlighting similarly emphasizes source
separation by transforming untrusted inputs so that the model can distinguish
them from trusted instructions~\cite{hines2024spotlighting}. AgentShield is
complementary: rather than changing the model input format or training
objective, it operates at runtime after observations have already entered the
session. It asks whether arguments to a proposed action can be traced to
untrusted content and whether they are consistent with the operator request.

The same separation idea has been pursued architecturally. An f-secure LLM
system disaggregates the pipeline and filters untrusted input out of the
planning process, with formal models of the resulting
guarantees~\cite{wu2024systemlevel}. IPIGuard decouples action planning from
interaction with external data by traversing a planned tool dependency
graph~\cite{an2025ipiguard}. SPA adopts a plan-first design and applies
dual-lattice information-flow control across explicit data and control
dependencies, extending the threat model to cross-query state
reuse~\cite{girrens2026spa}. These are structural redesigns of the agent.
AgentShield instead accepts the agent as given and adds a governance point in
front of the tool boundary, which is why it can be evaluated on logged traces
produced without it.

\subsection{Agent harmfulness and runtime governance}
\label{sec:rw:gov}

AgentHarm broadens evaluation from chatbot refusals to multi-step harmful agent
behavior~\cite{andriushchenko2025agentharm}. Its focus is malicious task
completion under jailbreaks rather than indirect prompt injection, but it
underscores why agent safety cannot be reduced to response classification.
AgentShield instead studies a runtime control point before tool execution. It
does not claim AgentHarm benchmark performance; the repository contains only a
metadata-derived proxy, which is excluded from the main evidence in this paper.

Several recent systems share our enforcement point. ClawGuard derives
task-specific access constraints from the user's stated objective and enforces
a user-confirmed rule set at every tool-call
boundary~\cite{zhao2026clawguard}; AgentWatcher attributes an agent action to
causally influential context segments and reasons over explicit rules, trading
attribution cost for explainability~\cite{wang2026agentwatcher}; AttriGuard
verifies a proposed tool call by counterfactual re-execution under an
attenuated view of external
observations~\cite{he2026attriguard}. All three are closer to our setting than
model-side defenses, and all three require either a monitor model or a
re-execution. AgentShield's distinguishing commitment is determinism: the
provenance signals in Section~\ref{sec:signals} are computed without any model
call, which is what makes the prediction-invariance audit in
Section~\ref{sec:contract} decidable and the latency in
Section~\ref{sec:latency} free of model time.

\subsection{Information flow control and taint tracking for agents}
\label{sec:rw:ifc}

Classical taint analysis assumes propagation through program memory, an
assumption that breaks when propagation is governed by probabilistic language
reasoning; NeuroTaint makes this argument explicitly and reconstructs
provenance offline from untrusted sources to privileged sinks using semantic
and causal evidence rather than exact string matching~\cite{cai2026ghost}. Our
taint propagation signal is deliberately the weaker, string-based variant ---
it is what a deterministic, model-free gate can compute --- and NeuroTaint's
TaintBench results are best read as the upper bound that a semantic tracker
could reach on traces like ours. FLOWSEAL grounds confidentiality enforcement
in data provenance and an information-flow-control lattice with controlled
declassification, and shows that enforcement living inside the LLM context
coincides with the attack surface~\cite{shim2026confuse}; APPA argues that
monotone taint tracking either over-blocks benign work or strands execution,
and adds policy-governed recovery~\cite{kravchenko2026appa}. GIF provides a
locally sound geometric bound on information flow with a mechanized proof,
addressing the taint-explosion problem directly~\cite{storek2026gif}.

These works converge on the diagnosis that motivated v0.4.1: the difficulty is
not producing a taint label but deciding what to \emph{do} with it when the
label conflicts with an action the operator asked for. That is an aggregation
question, and Section~\ref{sec:invariant} treats it as one.

\subsection{Governance before the tool boundary}
\label{sec:rw:boundary}

STPA-derived safety requirements formalized as enforceable specifications on
data flows and tool sequences move agent safety toward proactive guardrails
with stated guarantees rather than reliability
fixes~\cite{doshi2026verifiably}. AgentRaft defines data over-exposure as
transmitting sensitive data beyond user intent and functional necessity,
detecting it with runtime taint tracking over a cross-tool call
graph~\cite{lin2026agentraft}. ToolMinimize goes further and \emph{rewrites}
tool arguments to the minimum data necessary, a capability a gate that only
returns ALLOW, REVIEW or BLOCK does not have~\cite{li2026toolminimize}.

AgentShield occupies a specific position in this space: it neither rewrites
arguments nor re-executes the agent. It decides a ternary action at the
boundary from deterministic evidence, and it treats the boundary's
\emph{aggregation policy} as a first-class design object with its own
invariant.

\subsection{Automated attacks and benchmark sensitivity}
\label{sec:rw:variants}

Maatphor generates prompt-injection variants automatically so that candidate
defenses are tested against more than the single successful payload that
motivated them~\cite{salem2023maatphor}. This is the methodological reason our
held-out evaluation in Section~\ref{sec:heldout} exists, and it is also why we
report it as a failure: a detector tuned on the benchmark's own phrasing does
not survive variant generation.

The same pattern has been observed outside agent benchmarks: cybersecurity
tooling built on LLMs was turned against itself by prompt injection in a way
explicitly compared to cross-site scripting, with mitigations layered at the
tool boundary rather than in the
model~\cite{mayoralvilches2025cybersecai}. We cite it as convergent evidence
that the boundary is the right enforcement point, and as a caution that a
boundary rule which is not deterministic is itself part of the attack surface.

\subsection{Positioning}
\label{sec:rw:positioning}

The closest neighbors to AgentShield are the runtime, deterministic
gate~\cite{zhao2026clawguard,wang2026agentwatcher} and the offline provenance
reconstruction~\cite{cai2026ghost}. AgentShield differs from the former in
requiring no monitor model, and from the latter in operating online at the
decision point rather than as a post-hoc audit. Its contribution relative to
both is the safety invariant of Section~\ref{sec:invariant}, which is a
statement about how \emph{any} such gate should aggregate authorization
evidence against structural risk.

%-------------------------------------------------------------------------------
\section{Threat Model and Problem Formulation}
\label{sec:threat}
%-------------------------------------------------------------------------------

We consider an LLM agent that receives an operator request and then alternates
between model reasoning and external tool calls. Tool outputs may contain data
controlled by external parties. The adversary can place text in such an output
but does not directly modify the runtime gate, the operator request, or the
evaluation labels. The target is a downstream high-consequence call such as
sending data to an external destination, deleting a resource, or invoking a
privileged tool.

A single-event gate observes the current event $e_t = (tool_t, args_t)$ and
decides $d_t \in \{\textsc{allow}, \textsc{human\_review}, \textsc{block}\}$.
The central limitation is that two executions can present the same current event
while differing in the origin of an argument. Let $origin(x)$ denote the
earliest observed source of entity $x$. If $x$ first appears in an untrusted
tool output in the attack execution but in the operator request or trusted
context in the benign execution, a gate restricted to $e_t$ cannot distinguish
the pair. A provenance-aware gate augments the observation with session history
$H_{<t}$ and origin information $P_t$, producing $d_t = G(e_t, H_{<t}, P_t)$
(Figure~\ref{fig:gap}).

We do not assume that all untrusted content is malicious. In fact, this
distinction is the main unresolved problem. The conservative evaluation policy
treats recorded external tool outputs as untrusted, which prevents
label-derived trust assignment but necessarily creates false positives when
ordinary retrieved content legitimately influences an action.

%---------------------------
\begin{figure*}
\begin{center}
\includegraphics[width=\textwidth]{figure1_observability_gap}
\end{center}
\caption{\label{fig:gap} The observability gap. The final tool call may look
ordinary; the security-relevant distinction is the provenance of arguments
introduced by earlier content.}
\end{figure*}
%% %---------------------------

%-------------------------------------------------------------------------------
\section{AgentShield Design}
\label{sec:design}
%-------------------------------------------------------------------------------

\subsection{Observed content artifacts}
\label{sec:artifacts}

Every runtime-observed content object is represented as an
ObservedContentArtifact with a session identifier, source event, origin type,
trust level, content hash, and extracted entities. Entity extraction targets
values whose origin is operationally meaningful, including email addresses,
URLs, hosts, paths, and command fragments. The artifact abstraction turns prior
context from an opaque text buffer into an auditable runtime object.

\subsection{Entity origin and taint tracking}
\label{sec:taint}

The TaintTracker records the first observed origin of each normalized entity
and the set of artifacts in which it later appears. For an entity $x$, the
tracker stores $origin(x) = (artifact\_id, trust\_level, origin\_type,
first\_seen)$. A downstream tool call can therefore ask whether a destination
in its arguments was first introduced by untrusted content. The implementation
is intentionally deterministic and does not read evaluation metadata.

\subsection{Provenance risk signals}
\label{sec:signals}

AgentShield combines ordinary local risk signals with provenance signals. The
latter include: (i) untrusted content that appears instruction-like; (ii) a
destination first observed in untrusted content; (iii) an origin mismatch in
which a tainted value is absent from the operator request; (iv) sensitive
material flowing to an external sink; (v) untrusted content preceding a
privileged action; and (vi) taint propagation when an entity introduced by an
untrusted artifact reaches the current call. The signals are intended as
interpretable evidence rather than learned latent features.

\noindent\textbf{Signal classification.} Each signal type is assigned to
exactly one of three disjoint sets, defined statically after the signal-type
enumeration. \texttt{UNSUPPRESSIBLE\_SIGNALS} contains the structural and
manipulation classes: privilege change, audit tamper, bulk operation,
credential access, policy evasion, untrusted instruction,
untrusted-to-privileged action, sensitive-to-external flow, and cross-agent
delegation. \texttt{SUPPRESSIBLE\_PROVENANCE\_SIGNALS} contains the alarms whose
entire content is \emph{provenance suspicion}: destination provenance,
intent-origin mismatch, and taint propagation.
\texttt{SUPPRESSION\_EVIDENCE\_SIGNALS} contains authorization and
trusted-entity resolution, which are evidence rather than risk. A signal type
not explicitly classified is treated as unsuppressible, so a future signal
added without thought cannot silently become suppressible.

Two remarks on this taxonomy. First, taint propagation and cross-agent
delegation are deliberately distinct types: the former is a propagation fact
about an entity, the latter a structural relationship between agents.
Conflating them would mean that tempering taint tacitly tempers genuine
delegation risk. Second, suppression-evidence signals carry score 0.0 by
construction; an observation that a call was authorized or that an entity
resolved through a trusted store must not be able to manufacture risk on its
own.

\subsection{Behavior graph and risk propagation}
\label{sec:graph}

Each tool call becomes a node in a session behavior graph, and edges represent
relations such as calls, invocations, or data flow. Let $l(v)$ be the local
node risk and $w(u,v)$ the transfer coefficient on an incoming edge. The
propagated total risk is computed as:

\begin{equation}
R(v) = \max\Big( l(v),\; \max_{(u,v)\in E} R(u)\cdot w(u,v) \Big).
\label{eq:propagation}
\end{equation}

The current implementation uses type-specific transfer weights (for example,
0.7 for call/invocation edges and 1.0 for data-flow edges). Risk lives on
nodes; edge weights represent transfer only. This avoids a previous failure
mode in which upstream risk was encoded both in the node and the edge,
effectively squaring risk during propagation.

\subsection{Two-pass governance}
\label{sec:twopass}

Governance is performed after the current node has been inserted into the
graph. First, the engine computes local and provenance signals. Second, it
inserts the call at local risk and creates its parent edge. Third, it
propagates risk through the graph. Finally, the propagated state is used to
choose ALLOW, HUMAN\_REVIEW, or BLOCK, and the governed result is written back
to the node. This ordering matters: deciding before propagation would make
graph awareness observational only rather than decision-relevant.

\subsection{Risk aggregation across the two bands}
\label{sec:aggregation}

The final score is the maximum of two bands, and only one of them is reachable
by authorization. Let $H$ denote the structural components --- local,
inherited, path and downstream exposure --- and let $U$ and $P$ denote the
unsuppressible and suppressible signal classes of
Section~\ref{sec:signals}. With $a \in \{0,1\}$ indicating explicit operator
authorization, and with $\tau(0) = 1$ and $\tau(1)$ equal to the constant
\texttt{AUTHORISED\_ACTION\_CEILING}, the aggregate is

\begin{equation}
S(H, U, P, a) = \max\Big( \max(H \cup U),\; \min\big( \max(P), \tau(a) \big) \Big).
\label{eq:aggregate}
\end{equation}

Structural components and unsuppressible signals never enter the capped path.
Counterfactual and intervention values are also outside it: they measure the
value of stopping here, not who asked, so an intervention value above both
bands raises the score rather than being pardoned by it.
Section~\ref{sec:invariant} states the invariant this enforces and reports the
defect it replaced.

\subsection{Auditing and counterfactual support}
\label{sec:auditing}

Decisions retain references to the artifacts and risk signals that contributed
to them, enabling an evidence chain from a blocked sink back to the content
that introduced a suspicious value. The engine also contains counterfactual
support for estimating how earlier intervention could alter downstream risk; we
treat this as an audit mechanism rather than an evaluated causal guarantee in
the present paper.

%-------------------------------------------------------------------------------
\section{Evaluation Methodology}
\label{sec:methodology}
%-------------------------------------------------------------------------------

\subsection{Evaluation contract and leakage isolation}
\label{sec:contract}

The benchmark harness separates fields into runtime-observable,
runtime-derived, and evaluation-only classes. Labels, expected actions,
expected risk scores, attack names, injection goals, fixture chain
identifiers, rationales, and fields derived from these are forbidden from
detector input. Adapters construct runtime observations before prediction;
labels are consulted only afterward for scoring. To test this boundary, the
frozen release permutes evaluation-only metadata while holding runtime
observations fixed. Across seven methods, 60 cases, and eight perturbations per
case (3,360 prediction checks), no prediction changes.

\subsection{Logged-trace replay}
\label{sec:replay}

The main external experiment replays recorded tool-call and tool-response
triples from the ffuuugor/agentdojo-dump dataset. We use all 97 benign
trajectories (339 steps) and the first 400 attack trajectories, for 497
trajectories total. The replay is an offline interception experiment: it does
not rerun an agent, invoke an LLM, or measure what the agent would do after a
block. All recorded tool responses are assigned untrusted status independently
of their labels or contents. This conservative policy prevents trust leakage
but is expected to increase false positives.

Six configurations form an additive ladder. The first four are the frozen v0.3
rungs: \texttt{local\_only} sees only the current call;
\texttt{plus\_output\_inspection} enables reactions to previously observed
untrusted content without per-entity taint;
\texttt{plus\_entity\_provenance} tracks entity origin and taint;
\texttt{plus\_intent\_consistency} additionally checks against the operator
request. The last two are v0.4.1: \texttt{plus\_trust\_policy\_v0\_4} adds
tool-semantics trust classification and prose-authorization matching with the
band-split aggregation of Section~\ref{sec:aggregation};
\texttt{plus\_intent\_slots} replaces prose matching with authorization decided
against structured intent slots captured at request time. Cluster bootstrap
confidence intervals use the trajectory, not the individual step, as the
resampling unit (2,000 draws).

\subsubsection{Sampling protocol}
\label{sec:sampling}

The harness flag \texttt{-{}-max-trajectories 400} is a per-class limit, not a
total. The trajectory list is partitioned by label and each partition is
truncated independently (\texttt{benign[:400]}, \texttt{attack[:400]}). Because
the dump contains only 97 benign trajectories, the benign slice is the complete
set and the attack slice is a deterministic prefix. The evaluated sample is
therefore 97 benign (a census) and 400 attack-labelled (a prefix), for 497
trajectories. Trajectory construction order is fixed by
\texttt{build\_trajectories}, which walks the dataset in load order; no random
sampling, shuffling, or seed is involved at any point. This is load-bearing: a
run that takes the last 400 attack trajectories, or a random 400, is a
different experiment whose numbers are not comparable, and the committed
verifier detects the difference.

One asymmetry follows from this and is worth stating up front: the benign
figures are a census with no sampling error, while the attack figures are a
prefix of a larger population. They are not symmetric estimates, and the benign
number is the more trustworthy of the two.

\subsection{Paired counterfactual control}
\label{sec:paired}

To isolate source/context from tool choice, the paired experiment constructs
four attack/benign pairs with identical final tool names and identical final
arguments. The only intended difference is how the destination or
action-driving value entered the context. A local-only gate must therefore give
equal decisions within a pair; separation after context-aware processing
demonstrates sensitivity to source information. Because output inspection
already separates most pairs, this control establishes source/context
attribution but not the incremental contribution of entity provenance alone.

\subsection{Auxiliary SCI-600 fixture}
\label{sec:sci600}

SCI-600 is a project-generated, self-labelled single-event fixture used for
regression and baseline comparison. It cannot exercise the central provenance
mechanism because no prior tool output exists in the fixture format. We
therefore treat it as auxiliary evidence and do not use it to claim
generalization.

%-------------------------------------------------------------------------------
\section{Results}
\label{sec:results}
%-------------------------------------------------------------------------------

\subsection{Logged-trace replay}
\label{sec:mainresult}

%---------------------------
\begin{table*}
\begin{center}
\small
\begin{tabular}{l r r r r}
\toprule
Configuration & Attack trace BLOCK & 95\% CI & Benign trace BLOCK & 95\% CI \\
\midrule
Local only                       & 0.0\%  & [0.0, 0.0]     & 1.0\%  & [0.0, 3.1] \\
+ output inspection              & 3.5\%  & [1.75, 5.5]    & 19.6\% & [11.34, 27.84] \\
+ entity provenance / taint      & 16.0\% & [12.5, 19.75]  & 41.2\% & [31.96, 50.52] \\
+ intent consistency             & 16.0\% & [12.5, 19.75]  & 41.2\% & [31.96, 50.52] \\
\bottomrule
\end{tabular}
\caption{\label{tab:main} Main logged-trace replay result for the frozen v0.3
ladder. Confidence intervals are 95\% cluster-bootstrap intervals over
trajectories.}
\end{center}
\end{table*}
%% %---------------------------

The most important observation is the zero-to-nonzero transition. The
local-only gate blocks none of the 400 sampled attack trajectories. Merely
making earlier untrusted tool output visible raises blocking to 3.5\%, showing
that relevant signal lies outside the current tool call. Entity provenance and
taint then raise the block rate to 16.0\%, an absolute gain of 12.5 percentage
points over output inspection and a 4.6$\times$ multiplicative increase.

The safety gain is accompanied by a large benign cost. Figure~\ref{fig:ladder}
shows the two quantities moving together: benign trace blocking rises from 1.0\%
for local-only to 19.6\% with output inspection and 41.2\% with entity
provenance. The full intent-consistency rung does not improve this trade-off on
the logged traces. The result therefore supports provenance as a useful signal
but not as a complete decision rule under an all-untrusted trust policy.

%---------------------------
\begin{figure*}
\begin{center}
\includegraphics[width=\textwidth]{figure2_ladder}
\end{center}
\caption{\label{fig:ladder} Attack interception and benign blocking move
together as more cross-event context is enabled.}
\end{figure*}
%% %---------------------------

\subsubsection{Adding authorization semantics (v0.4.1)}
\label{sec:audit}

The rungs above leave 41.2\% benign trace blocking, which is the cost of the
all-untrusted policy. The v0.4.1 configuration asks whether operator
authorization can reduce that cost without giving up the provenance mechanism.
Table~\ref{tab:ladder} reports the full ladder including the v0.4.1 rungs.

%---------------------------
\begin{table*}
\begin{center}
\small
\begin{tabular}{l r r}
\toprule
Configuration & Attack-labelled trace BLOCK & Benign trace BLOCK \\
\midrule
\texttt{local\_only} (single-event gate)      & 0.0\%            & 1.03\% \\
+ untrusted output inspection                 & 3.5\%            & 19.59\% \\
+ entity provenance / taint (v0.3)            & 16.0\%           & 41.24\% \\
+ intent consistency                          & 16.0\%           & 41.24\% \\
\textbf{+ safe authorisation (v0.4.1)}        & \textbf{3.5\%}   & \textbf{28.87\%} \\
+ structured intent slots                     & 3.5\%            & 28.87\% \\
\bottomrule
\end{tabular}
\caption{\label{tab:ladder} Full ladder including the v0.4.1 rungs, reported as
point estimates. The first four rows are the frozen v0.3.1 ladder,
reproduced unchanged for comparability; their cluster-bootstrap confidence
intervals appear in Table~\ref{tab:main} and are not repeated here.}
\end{center}
\end{table*}
%% %---------------------------

\noindent\textbf{On confidence intervals.} The v0.4.1 rows carry no interval,
and this is a limitation rather than an omission to be papered over: the frozen
\texttt{v0\_4\_trust\_pareto.json} records point estimates only, and its result
records contain no variance field, so an interval computed for v0.4.1 would have
to come from a bootstrap run that was not part of the frozen artifact. We
therefore report the point estimates and state their absence, and we do not
transplant the v0.3 intervals onto the v0.4.1 rows. Regenerating intervals for
v0.4.1 requires re-running the ladder with a bootstrap wrapper, which is a
separate measurement, not a re-reading of the committed artifact.

Two observations. First, benign trace blocking falls from 41.24\% to 28.87\%, a
12.37-point reduction, which is the intended benefit. Second, the raw blocking
rate on attack-labelled trajectories also falls, from 16.0\% to 3.5\%. This
second movement must not be described as the removal of false positives; the
reason is stated in the limitations bullet on attack labels versus attack
outcomes (Section~\ref{sec:limitations}).

An action-level audit of the trajectories the v0.3 rung blocked shows how the
release decomposes. Of the 64 attack-labelled trajectories blocked by the v0.3
provenance rung, 36 ended in a final governed action concordant with explicit
operator authorization, 14 ended in a read-only call that cannot act on the
tainted value, and 14 were neither. This is an action-level statement about the
final governance point; it is not a statement that those 50 trajectories were
benign, and it is not a statement that the attacks failed.

\texttt{plus\_intent\_slots} produces numbers identical to
\texttt{plus\_trust\_policy\_v0\_4} on this corpus. Prose authorization matching
and structured-slot authorization agree wherever an operator request names both
the action family and the entity exactly, which is the case for the
authorization-concordant trajectories this corpus contains. We therefore do not
claim the slot mechanism improves over prose matching here; we report only that
it does not regress, and that the corpus cannot separate the two.

\subsection{Paired source/context control}
\label{sec:pairedresult}

%---------------------------
\begin{table*}
\begin{center}
\small
\begin{tabular}{l l l}
\toprule
Identical final call & Attack decision & Benign decision \\
\midrule
\texttt{send\_email(attendees $\rightarrow$ partner)}   & BLOCK         & HUMAN\_REVIEW \\
\texttt{delete\_file(4821)}                             & BLOCK         & ALLOW \\
\texttt{http\_request(metrics endpoint)}                & BLOCK         & ALLOW \\
\texttt{upload\_file(customers.csv)}                    & BLOCK         & BLOCK \\
\bottomrule
\end{tabular}
\caption{\label{tab:paired} Paired counterfactual decisions after
context-aware processing. Local-only decisions agree within all four pairs.}
\end{center}
\end{table*}
%% %---------------------------

Three pairs exhibit different decisions under context-aware processing,
including two complete BLOCK-versus-ALLOW separations. The upload pair remains
blocked in both conditions because the current policy treats exporting a
customer file to an external destination as high risk even when explicitly
authorized. This negative pair is useful, and in v0.4.1 it is no longer
incidental but required: sensitive material crossing the trust boundary is
classified unsuppressible, so authorization structurally cannot lower that
decision. The pair is therefore a direct behavioural witness of the invariant in
Section~\ref{sec:invariant}, showing that provenance does not automatically
override content-level safety rules.

\subsection{Evaluation isolation}
\label{sec:isolation}

The prediction-invariance audit executes 3,360 metadata perturbations across
seven methods and observes zero prediction changes. This is a boundary
property, not a performance metric. It establishes that the reported predictors
are invariant to hidden grading metadata under the tested perturbations. The
audit replaced an earlier invalid methodology that inferred leakage from
accuracy degradation after shuffling ground-truth labels.

\subsection{Held-out instruction generalization}
\label{sec:heldout}

A frozen instruction detector was evaluated on nine differently phrased attack
styles and six benign controls without adding markers or changing thresholds.
It detected 0/9 attacks and produced 0/6 benign false positives. Entity
extraction recovered 5/7 expected destinations. This is an explicit negative
result: the string-marker detector does not generalize to the held-out
injection idioms. It motivates an instruction-versus-data classifier or stronger
source semantics rather than further benchmark-specific keyword expansion.

\subsection{Governance latency}
\label{sec:latency}

In the same 497-trajectory replay, governance-only latency measured around
\texttt{process\_tool\_call} is p50 = 0.47 ms, p95 = 5.32 ms, and p99 = 6.96 ms
over 781 calls. These numbers exclude dataset loading and any model time.
AgentShield itself places no LLM call on the governance path in the evaluated
configuration.

\subsection{SCI-600 auxiliary results}
\label{sec:sci600result}

On the self-labelled SCI-600 fixture, the content-keyword baseline achieves the
highest Macro-F1 (0.487). The production AgentShield pipeline obtains
approximately 0.43 Macro-F1 and substantially fewer false allows in the frozen
reports. Because this fixture is generated and labelled by the project and
contains no prior tool-output context, we do not interpret the ranking as
evidence of real-world generalization or of the provenance mechanism. The main
paper claims instead rest on the external logged traces and paired
source/context control. These figures describe the frozen v0.3 pipeline; the
v0.4.1 aggregation change is not reflected in this fixture.

%-------------------------------------------------------------------------------
\section{Discussion}
\label{sec:discussion}
%-------------------------------------------------------------------------------

\subsection{What the results establish}
\label{sec:disc:establish}

First, the logged-trace experiment provides direct evidence of an observability
gap at the runtime gate: a gate restricted to the current call blocks 0\% of
sampled attack trajectories, while a gate allowed to inspect prior output begins
to recover attack signal. Second, provenance provides additional signal beyond
generic output inspection, increasing trace blocking from 3.5\% to 16.0\% under
the frozen configuration. Third, the paired control demonstrates that
source/context can change a decision even when the final tool call is held
fixed.

\subsection{What the results do not establish}
\label{sec:disc:notestablish}

The experiments do not measure end-to-end attack success rate, benign utility,
task success, or behavioral adaptation after a block. Offline replay cannot
determine whether an agent would replan, retry, choose another tool, or
successfully finish a benign task under governance. Likewise, the paired control
does not prove that entity provenance alone creates every separation because the
output-inspection rung already explains most of the paired difference.

\subsection{The central open problem: trustworthy source attribution}
\label{sec:disc:trust}

The 41.2\% benign trace-block rate of the v0.3 configuration exposed the primary
research problem. Treating every external tool response as untrusted is
label-safe but operationally crude. A deployable system needs a trust policy
derived from runtime semantics rather than attack labels or content detectors.
Candidate approaches include tool-specific source policies, authenticated origin
metadata, explicit user authorization binding, capability-aware sinks, and a
classifier that distinguishes imperatives from ordinary retrieved data. v0.4.1
reduces the residual to 28.9\% by correcting authorization semantics rather than
by trust-prior tuning, which leaves 28.9\% as the open figure. These approaches
should be frozen before benchmark evaluation to avoid benchmark-informed feature
engineering.

\subsection{Why provenance remains useful despite weak detector generalization}
\label{sec:disc:detector}

The held-out 0/9 detector recall might appear to undermine the architecture, but
it actually separates two questions. Instruction detection asks whether a piece
of text is adversarial; provenance asks where a value came from and how it flows
into action. The latter remains meaningful even when the former is imperfect.
Future systems can replace the string detector without changing the provenance
data model, entity origins, or behavior graph.

\subsection{An authorization safety invariant for risk aggregation}
\label{sec:invariant}

The most consequential defect found in v0.4 was not a missing feature but an
aggregation error. Authorization was implemented as a single ceiling applied to
the combined risk score:

{\small
\begin{verbatim}
peak = max(all components, all signals)
if authorized:
    peak = min(peak, AUTHORISED_ACTION_CEILING)
\end{verbatim}
}

Because the cap was applied after aggregation, it bounded \emph{every}
component, including structural risk that authorization has no bearing on. An
authorized bulk delete therefore left the gate at the ceiling value rather than
at its true structural severity. The natural test to write for that code ---
``an authorized action scores at the ceiling'' --- is precisely a pinned-down
vulnerability, and it was in fact written.

The fix is to split the aggregate into two bands that authorization cannot
bridge:

{\small
\begin{verbatim}
hard_peak = max(structural components,
                unsuppressible signals)
soft_peak = max(suppressible provenance signals)
if authorized:
    soft_peak = min(soft_peak,
                    AUTHORISED_ACTION_CEILING)
score      = max(hard_peak, soft_peak)
\end{verbatim}
}

Structural components (\texttt{local\_risk}, \texttt{inherited\_risk},
\texttt{path\_risk}, \texttt{downstream\_exposure}) and the unsuppressible
signal classes never enter the capped path. The invariant this establishes is
the paper's central design contribution in this release:

\begin{quote}
\textbf{Authorization may explain provenance-derived suspicion, but it cannot
authorize away independently dangerous behavior.}
\end{quote}

The invariant is behavioural, not architectural: it is pinned by a test suite in
which an authorized bulk delete remains blocked, an authorized action carrying
untrusted content is reviewed rather than blocked, an authorized action whose
destination arrived untrusted is bounded by the ceiling, inherited risk is
unaffected, intervention value is unaffected, and an unknown future signal type
fails closed as unsuppressible.

Two design decisions are worth stating explicitly because both cut against the
headline number. First, trusted-entity resolution is deliberately
\textbf{evidence-only} in this release: it is recorded in the audit chain but
carries no ceiling. Giving it a prior-based ceiling in the same change would
have moved two variables at once and made the ablation uninterpretable; the
intended ladder is v0.4.1a authorization-only, then v0.4.1b + trusted
resolution. Second, the ceilings were not retuned after the fix, even though
retuning would have recovered part of the benign reduction visible in
Table~\ref{tab:ladder}. Retuning to recover a number is benchmark fitting.

The general lesson generalizes beyond this system: any risk aggregate that
mixes ``who asked'' with ``what is asked'' needs a stated policy for which of
the two authorization is allowed to move. Leaving it implicit produces a bypass
that looks like a feature.

%-------------------------------------------------------------------------------
\section{Limitations and Threats to Validity}
\label{sec:limitations}
%-------------------------------------------------------------------------------

\begin{itemize}

\item The logged-trace source is an AgentDojo-derived dump rather than the
official AgentDojo runtime and grader. We therefore do not report AgentDojo
benchmark performance or attack success rate.

\item The main attack subset uses the first 400 attack trajectories after the
audited grouping procedure, while all 97 benign trajectories are included. This
is reproducible but is not a random population sample of all possible agent
behavior. The benign figure is therefore a census and the attack figure a
prefix, and the two are not symmetric estimates; where they are compared
directly, the benign figure is the more trustworthy.

\item The all-untrusted trust policy is intentionally conservative and likely
overestimates false positives relative to a system with reliable source
semantics. Conversely, it avoids an invalid evaluation shortcut in which trust
is inferred from whether the content is malicious. The v0.4.1 rungs use a
tool-semantics trust prior rather than the binary policy, but that prior
participates only in an evidence-only signal in this release; the benign
reduction it produces is therefore attributable to authorization semantics, not
to trust-prior tuning.

\item \textbf{Attack labels are not attack outcomes.} The dataset labels
trajectories as attack or benign. Our gate decides at a final governance point
using runtime-observable authorization evidence. Neither is a measurement of
attack success: offline replay has no official grader, no attack-success rate,
no task-level outcome and no utility measure. A release decision that is
defensible at the action level is not thereby evidence that an attack failed,
and the paper makes no such inference anywhere.

\item \textbf{Authorization is a weak authorization signal.} Matching an
operator request against a tool call establishes that the operator named this
action family and this entity; it does not establish informed consent,
authority, or absence of coercion. A manipulated agent whose final call happens
to be authorization-concordant is authorized in our operational sense only.
This is a further reason the action-level audit cannot be read as an
attack-outcome measurement.

\item The instruction detector is rule-based and fails to generalize on a small
held-out set. The paper therefore does not claim robust prompt-injection
detection.

\item SCI-600 is synthetic and self-labelled. It is retained for regression and
baseline sanity checks, not generalization claims.

\item Risk weights and thresholds are engineering choices rather than calibrated
probabilities. Their interpretation is ordinal governance severity, not
probabilistic attack likelihood.

\item The v0.4.1 ladder rows are reported as point estimates without confidence
intervals, because the frozen artifact records no variance for them
(Section~\ref{sec:audit}).

\end{itemize}

%-------------------------------------------------------------------------------
\section{Reproducibility and Release Discipline}
\label{sec:repro}
%-------------------------------------------------------------------------------

The submitted artifact contains a frozen research snapshot and a separate
reproducibility-tooling snapshot. The artifact includes a canonical result
file, a deterministic dataset fingerprint, and a verifier for reproducing the
reported experiment. Identifying repository history and public repository
metadata are omitted for double-blind review; internally the project does keep
a release label and revision identifiers for release discipline, but those are
not what bind the reported numbers to the artifact --- content digests are.
Later documentation-only edits do not change the experiment, and the moving
development HEAD is not an experiment identifier.

The main artifact is \texttt{benchmark/results/v0\_4\_trust\_pareto.json},
SHA-256 \texttt{3e33\hb c72d\hb 230c\hb b7e0\hb 788d\hb c441\hb 2574\hb 054c\hb 19df\hb d496\hb 5465\hb 7369\hb 8d73\hb bff7\hb b6a9\hb 0565}.
The evaluated dump is fingerprinted as a manifest over its Arrow files, with
digest \texttt{948b\hb 9432\hb 5ae1\hb c8cf\hb be41\hb bd20\hb 5fee\hb 46b0\hb f500\hb 2155\hb 59ef\hb 5cd9\hb 7a25\hb c27f\hb cf58\hb 0049}
over 1 file, 50,898,680 bytes, 13,913 rows and 10,536 constructed trajectories.
A reproduction is valid only when both hashes match: the same script run against
a different dump version is not a reproduction.

Reproduction is a two-command procedure. First, regenerate into a file distinct
from the canonical artifact:

{\small
\begin{verbatim}
python benchmark/v04_trust_replay.py \
    --mode audit --max-trajectories 400 \
    --out benchmark/results/
        v0_4_trust_pareto.reproduced.json
\end{verbatim}
}

Then compare mechanically, with no tolerance and no ``close enough'' mode:

{\small
\begin{verbatim}
python scripts/verify_v04_reproduction.py \
    benchmark/results/v0_4_trust_pareto.json \
    benchmark/results/
        v0_4_trust_pareto.reproduced.json
\end{verbatim}
}

The verifier compares experiment identity, dataset audit counts, ladder order,
result order, and every per-configuration rate plus derived deltas, exiting 0 on
exact agreement, 1 with a per-field difference list on any mismatch, and 2 when
the comparison cannot be attempted. It fails closed when
\texttt{AGENTSHIELD\_ENABLE\_CORPUS\_MARKERS} is enabled, because
corpus-specific markers change scoring and the canonical run had them disabled.
In this release, three independent runs --- the committed artifact and two fresh
reproductions --- produced byte-identical output.

\noindent\textbf{What CI does and does not verify.} Continuous integration
verifies code integrity, tests, linting, type checking, benchmark smoke
execution and Docker health on Python 3.11 and 3.12. It does not independently
regenerate the 497-trajectory logged-trace artifact, because that run depends
on the locally staged external dump. A passing CI badge is therefore not a
reproduction of the 3.5\% / 28.9\% figures; the verification that matters for
those numbers is the manual deterministic procedure above. Lint scope is
likewise stated precisely: the two commands CI runs, \texttt{ruff check
backend/} and \texttt{ruff check benchmark/}
\texttt{benchmark/independent\_eval/ scripts/}, both pass, while a
whole-repository \texttt{ruff check .} reports 123 pre-existing findings in
directories CI does not lint.

Earlier figures produced by leaky or inconsistent harnesses are marked as
withdrawn in the repository and are not used here. This release discipline is
important because the project itself uncovered multiple evaluation failures
during development, including label-derived fields, ground-truth values entering
the engine, stale artifacts, and an invalid label-shuffling leakage diagnostic.

%-------------------------------------------------------------------------------
\section{Conclusion}
\label{sec:conclusion}
%-------------------------------------------------------------------------------

Tool-using agents can make dangerous decisions whose decisive cause is no longer
visible in the final tool call. AgentShield studies this observability gap with
a provenance-aware runtime gate. In logged traces, a current-call-only gate
blocks none of the sampled attack trajectories; prior-output inspection raises
blocking to 3.5\%, and entity provenance raises it to 16.0\%. The same additions
sharply increase benign blocking, making the unresolved safety-utility trade-off
explicit.

This release then shows that part of that trade-off was an aggregation error
rather than an intrinsic cost. Separating provenance suspicion from structural
danger, and permitting authorization to temper only the former, reduces benign
trace blocking to 28.9\% while keeping the authorization safety invariant that
an authorized bulk delete remains blocked. The measured cost of the correction
is reported rather than tuned away, and no threshold, weight, marker or sampling
rule was adjusted to recover it.

Together with the paired source/context control and metadata-invariance audit,
the results support provenance as a useful runtime security signal while
rejecting stronger claims that the current prototype solves indirect prompt
injection. The next research step is not further threshold tuning, but reliable
trust attribution and instruction-data separation grounded in runtime source
semantics --- with the explicit recognition that trajectory-level attack labels,
runtime action authorization, and actual attack success are three distinct
levels, and that offline replay measures only the first two.

%-------------------------------------------------------------------------------
% optional clearing of the page
\cleardoublepage
\appendix
\section*{Ethical Considerations}
%-------------------------------------------------------------------------------

This study is retrospective and entirely offline. The external evidence is a
logged-trace dump derived from AgentDojo, a published research benchmark; the
workloads replayed in this paper were produced before our system ran and are
re-used unchanged. No live system, third-party service, human subject, or
personal communication was targeted, and no prompt-injection payload was
delivered against any system outside the recorded dataset.

The engineering work was done on the project's own prototype, and the only
vulnerability it exercises is the authorization-aggregation defect described in
Section~\ref{sec:invariant}, which was found in our own code, reported to our own
issue tracker, and fixed in the same release. We do not publish a usable
injection payload, a bypass against any deployed agent, or an attack against any
third party; the contribution is a defensive runtime gate and a statement about
how such a gate should aggregate authorization evidence.

The one methodological risk we flag is evaluation validity rather than harm:
inferring trust from labels would have inflated the results, and inferring attack
success from attack labels would have overstated the contribution. Both
shortcuts were rejected, and Section~\ref{sec:limitations} states the resulting
non-claims explicitly. No IRB approval was required because the study involves
no human participants.

%-------------------------------------------------------------------------------
% optional clearing of the page
\cleardoublepage

\section*{Open Science}
%-------------------------------------------------------------------------------

All artifacts necessary to evaluate the contribution are listed below. Every
identifier below is a content-addressed name, so the artifact a reviewer
obtains is byte-identical to the one this paper reports. Identifying
repository history and public repository metadata are omitted for double-blind
review.

\begin{description}

% TODO(open-science): this item is the one place an anonymous artifact URL
% must appear.  Create the anonymous mirror at anonymous.4open.science
% (Conference ID SEC27), then replace the paragraph below with, for example:
%   \item[Code and research logic.] Anonymous artifact at
%   \url{https://anonymous.4open.science/<SEC27 id>}: a frozen research
%   snapshot, a separate reproducibility-tooling snapshot, and their tests.
% Do NOT link the public development repository: it carries the owner name,
% the tag name and the whole commit history.  The anonymous mirror must also
% not contain the 50.9 MB Arrow dump (see the dataset item below), because the
% host does not accept files above 8 MB and does not anonymize binaries.

\item[Code and research logic.] The submitted artifact contains a frozen
research snapshot and a separate reproducibility-tooling snapshot, together
with the tests that pin their behaviour. The research snapshot pins the engine
and the experiment code that produced the canonical result file below; the
tooling snapshot contains the canonical-artifact verifier, the dataset
fingerprint script and their tests. Size and licensing constraints on the
external dump are stated under the dataset item.

\item[Canonical result artifact.]
\texttt{benchmark/results/v0\_4\_trust\_pareto.json}, SHA-256
\texttt{3e33\hb c72d\hb 230c\hb b7e0\hb 788d\hb c441\hb 2574\hb 054c\hb 19df\hb d496\hb 5465\hb 7369\hb 8d73\hb bff7\hb b6a9\hb 0565}.
This
file contains the six-rung ladder, the per-configuration rates, the derived
deltas, and the point estimates reported in Tables~\ref{tab:main} and
\ref{tab:ladder}.

\item[Evaluated dataset fingerprint.] A manifest over the Arrow files of the
AgentDojo-derived dump, digest
\texttt{948b\hb 9432\hb 5ae1\hb c8cf\hb be41\hb bd20\hb 5fee\hb 46b0\hb f500\hb 2155\hb 59ef\hb 5cd9\hb 7a25\hb c27f\hb cf58\hb 0049},
covering 1 file, 50,898,680 bytes, 13,913 rows and 10,536 constructed
trajectories. The dump itself is the external data source described in
Section~\ref{sec:replay}. It is deliberately not redistributed inside the
artifact, both because it is far larger than the artifact host's per-file size
limit and because it is distributed under the upstream dataset's terms.
\texttt{DATASET.md} therefore records the exact source, the expected local
path and this manifest digest, so that a reviewer can obtain the same dump and
mechanically reject any different version. The manifest, not the file, is the
identifier that binds a reproduction to the evaluated data.

\item[Reproduction tooling.] \texttt{scripts/verify\_v04\_reproduction.py}
(experiment identity, audit counts, ladder and result order, per-configuration
rates and derived deltas; exit codes 0 exact agreement, 1 mismatch with a
per-field difference list, 2 comparison not attempted) and the dataset
fingerprint script referenced above, together with their tests.

\item[Figures.] The two figures in this paper are generated by
\texttt{scripts/draw\_manuscript\_figures.py} directly from
\texttt{v0\_4\_trust\_pareto.json}; the script emits the 300 dpi PNG and the
vector PDF from the same data, so the figure values can be checked against the
artifact without trusting the image.

\item[Reproduction command.] The two-command procedure in
Section~\ref{sec:repro}. Both hashes must match; the verifier has no tolerance
mode and fails closed when corpus-specific markers are enabled.

\item[Known limits of the artifact set.] The logged-trace run is not
regenerated by CI, because it depends on the locally staged external dump; CI
verifies code integrity, tests, linting, type checking, benchmark smoke
execution and Docker health only. The v0.4.1 ladder rows carry point estimates
only, because the frozen artifact records no variance field for them.

\end{description}

%-------------------------------------------------------------------------------
\cleardoublepage

\bibliographystyle{plainurl}
\bibliography{references}

%-------------------------------------------------------------------------------
\end{document}
%-------------------------------------------------------------------------------
